\chapter{Optimiser, comparer et approcher}
\section*{Exercice 11.1 -- Boîte sans couvercle}
Dans un carré de côté $6$, on enlève aux coins des carrés de côté $x$ et on replie. Maximiser le volume.

\section*{Exercice 11.2 -- Comparer deux expressions}
Démontrer $x+1/x\ge2$ pour $x>0$, algébriquement puis par dérivation.

\section*{Exercice 11.3 -- Erreur d'une tangente}
Pour $f(x)=x^3$, approcher $f(1{,}02)$ par sa tangente en $1$ et borner l'erreur avec $|f''|\le6{,}12$.

\section*{Exercice 11.4 -- Newton pour une racine}
À partir de $u_0=2$, calculer deux pas de Newton pour $\sqrt2$ et fournir un encadrement.

\section*{Exercice 11.5 -- Échec de Newton}
Appliquer Newton à $f(x)=x^3-2x+2$ à partir de $0$. Que révèle cet exemple ?

\section*{Exercice 11.6 -- Maximum sur une grille}
Une fonction vérifie $|f'|\le5$ sur $[0,1]$. Une grille contient les extrémités et a un pas de $0{,}02$. Borner l'écart entre le maximum vrai et le maximum lu sur la grille.
